\documentclass[11pt]{article}
\usepackage[a4paper, top=18mm, bottom=18mm, left=20mm, right=20mm]{geometry}
\usepackage[T2A]{fontenc}
\usepackage[utf8]{inputenc}
\usepackage[ukrainian]{babel}
\usepackage{graphicx}
\usepackage[export]{adjustbox}
\usepackage{amsmath}
\usepackage{amssymb}
\graphicspath{{/app/Quiz/Scripts/2023/ProgAll/15BinTree/}}
\setlength{\parindent}{0pt}
\setlength{\parskip}{6pt}
\emergencystretch=3em
\pagestyle{empty}
\begin{document}
%begin
{\centering\Large\bfseries Контрольна робота\par}
\medskip
\noindent\rule{\linewidth}{0.5pt}
\smallskip
\noindent\textbf{Варіант \No\,1}\hfill \textit{ПІБ:}\,\underline{\hspace{60mm}}
\vspace{4mm}

%beginitem
1. %goodpath:Scripts/2018/Mod2018_1/03-good.txt
\textbf{Відмітити все, що є коректним оператором С++:}
( 1 ) {a=5;}

( 2 ) char c='\n';

( 3 ) int f(int a,9);

( 4 ) if a<b a=0;
\par\vspace{5mm}
%enditem
%beginitem
2. \textbf{Що буде виведено за виконання фрагменту коду?}

%code
4**-0.5*True
%code
\par\vspace{5mm}
%enditem
%beginitem
3. \textbf{Що буде виведено за виконання фрагменту коду?}

%code
print(6 != 2 > 4 < False)
%code
\par\vspace{5mm}
%enditem
%beginitem
4. %goodpath:Scripts/2019/Mod2019_1/Lex/01-good.txt
\textbf{Відмітити все, що є однією правильною лексемою:}
( 1 ) 0.5

( 2 ) or

( 3 ) -3J

( 4 ) "abc"""
\par\vspace{5mm}
%enditem
%beginitem
5. \textbf{Що буде виведено за виконання фрагменту коду?}

a,b,c=4,0,2
def g(b):
global c    a-=2
    b=3
    c=1
    return a+b+c

a,b,c=5,1,6
print(g(b),a,b,c)
\par\vspace{5mm}
%enditem
%beginitem
6. \textbf{Що буде виведено за виконання фрагменту коду?}

%code
cout << (9 / 9 / 9 / 9);
%code
\par\vspace{5mm}
%enditem
%beginitem
7. 

\begin{center}
	\adjustimage{max size={0.8\linewidth}{0.8\paperheight}}
	{../15BinTree/trees_exp/40.png}
\end{center}
Указати послідовність міток вершин дерева, зображеного на рис., що утвориться в результаті обходу дерева в оберненому порядку. Мітки записувати через пробіл.\par Алгоритм починає роботу з кореня (на рис. це верхня вершина).
%answer:40 оберненому
\par\vspace{5mm}
%enditem
%beginitem
8. %goodpath:Scripts/2019/Mod2019_1/Lex/03-good.txt
\textbf{Відмітити все, що є коректним оператором С++:}
( 1 ) x=y=z

( 2 ) ++n

( 3 ) True=1

( 4 ) n=1; n++
\par\vspace{5mm}
%enditem
%beginitem
9. \textbf{Що буде виведено за виконання фрагменту коду?}
%code
a, b, c = 7, 6, 8
def g(a, b, c=9):
    print(a, b, c, end=" ")

g(a=4, 0, b=2)
print(a, b, c)
%code
\par\vspace{5mm}
%enditem
%beginitem
10. \textbf{Що буде виведено за виконання фрагменту коду?}

print('R', end=' ')
e=['0','1','2','3','4']
e.extend([1,2])
print(e)

\par\vspace{5mm}
%enditem










%end
\newpage

\newpage

%begin
{\centering\Large\bfseries Контрольна робота\par}
\medskip
\noindent\rule{\linewidth}{0.5pt}
\smallskip
\noindent\textbf{Варіант \No\,2}\hfill \textit{ПІБ:}\,\underline{\hspace{60mm}}
\vspace{4mm}

%beginitem
1. \textbf{Що буде виведено за виконання фрагменту коду?}

a,b,c=4,0,2
def g(b):
global c    a-=2
    b=3
    c=1
    return a+b+c

a,b,c=5,1,6
print(g(b),a,b,c)
\par\vspace{5mm}
%enditem
%beginitem
2. \textbf{Що буде виведено за виконання фрагменту коду?}

%code
print('R', end=' ')
e = ('a','b','c','d','e',0,1,2)
print(e[-8])
%code

\par\vspace{5mm}
%enditem
%beginitem
3. %goodpath:Scripts/2018/Mod2018_1/03-good.txt
\textbf{Відмітити все, що є коректним оператором С++:}
( 1 ) {a=5;}

( 2 ) char c='\n';

( 3 ) int f(int a,9);

( 4 ) if a<b a=0;
\par\vspace{5mm}
%enditem
%beginitem
4. \textbf{Що буде виведено за виконання фрагменту коду?}

%code
cout << (!6==2 or !4< false or 6.0>2.0);
%code
\par\vspace{5mm}
%enditem
%beginitem
5. \textbf{Що буде виведено за виконання фрагменту коду?}
%code

a,b,c=7,6,8
def g(a,b,c=9):
    print(a,b,c,end="")

g(a=4,0,b=2)
print(a,b,c)

%code
\par\vspace{5mm}
%enditem
%beginitem
6. \textbf{Що буде виведено за виконання фрагменту коду?}

x,y,z=4,0,2
y,x,z,z=y,7,z,x
print(x,y,z)
\par\vspace{5mm}
%enditem
%beginitem
7. %goodpath:Scripts/2019/Mod2019_1/Lex/03-good.txt
\textbf{Відмітити все, що є коректним оператором С++:}
( 1 ) x=y=z

( 2 ) ++n

( 3 ) True=1

( 4 ) n=1; n++
\par\vspace{5mm}
%enditem
%beginitem
8. %goodpath:Scripts/2018/Mod2018_1/01-good.txt
\textbf{Відмітити все, що є однією правильною лексемою:}
( 1 ) False

( 2 ) fop

( 3 ) <>

( 4 ) '123'
\par\vspace{5mm}
%enditem
%beginitem
9. \textbf{Що буде виведено за виконання фрагменту коду?}

%code
print(4**-0.5*True)
%code
\par\vspace{5mm}
%enditem
%beginitem
10. \textbf{Що буде виведено за виконання фрагменту коду?}

%code
#include <iostream>
using namespace std;

int g(int b){
    int y = 50;
    if (b) 
        y = 4;
    else if (b == 0)
         return 0;
    else
         y = 2;
    return y;
}

int main(){
    cout << g(-4);
    return 0;
}
%code

\par\vspace{5mm}
%enditem










%end
\newpage

\newpage

%begin
{\centering\Large\bfseries Контрольна робота\par}
\medskip
\noindent\rule{\linewidth}{0.5pt}
\smallskip
\noindent\textbf{Варіант \No\,3}\hfill \textit{ПІБ:}\,\underline{\hspace{60mm}}
\vspace{4mm}

%beginitem
1. \textbf{Що буде виведено за виконання фрагменту коду?}

%code
cout << 9 / 9 / 9 / 9;
%code
\par\vspace{5mm}
%enditem
%beginitem
2. %goodpath:Scripts/2019/Mod2019_1/Lex/01-good.txt
\textbf{Відмітити все, що є однією правильною лексемою:}
( 1 ) 0.5

( 2 ) or

( 3 ) -3J

( 4 ) "abc"""
\par\vspace{5mm}
%enditem
%beginitem
3. \textbf{Що буде виведено за виконання фрагменту коду?}
try:
    res = 4**-0.5*True
    if isinstance(res, bool):
        print(f'bool;{res}')
    elif isinstance(res, int):
        print(f'int;{res}')
    elif isinstance(res, float):
        print(f'float;{res:.2f}')
    else:
        print('???')
except:
    print('Z')
\par\vspace{5mm}
%enditem
%beginitem
4. \textbf{Що буде виведено за виконання фрагменту коду?}
try:
    res = 8//8/8*8
    if isinstance(res, bool):
        print(f'bool;{res}')
    elif isinstance(res, int):
        print(f'int;{res}')
    elif isinstance(res, float):
        print(f'float;{res:.2f}')
    else:
        print('???')
except:
    print('Z')
\par\vspace{5mm}
%enditem
%beginitem
5. 
Вкажіть значення та тип виразу:
"6.0j"==6
\par\vspace{5mm}
%enditem
%beginitem
6. %goodpath:Scripts/2019/Mod2019_1/Lex/03-good.txt
\textbf{Відмітити все, що є коректним оператором С++:}
( 1 ) x=y=z

( 2 ) ++n

( 3 ) True=1

( 4 ) n=1; n++
\par\vspace{5mm}
%enditem
%beginitem
7. \textbf{Що буде виведено за виконання фрагменту коду?}
try:
    for e in range(7, 7, -2):
        if e<=7:
            continue
        print(e, end=';');e = -5
        if e<=7:
            break
    else:
        print(e,end=';')
    print(e)
except:
    print('Z')
\par\vspace{5mm}
%enditem
%beginitem
8. \textbf{Що буде виведено за виконання фрагменту коду?}
%code
#include <iostream>
using namespace std;

int f(int &x, int &y){
    x = 6;
    y-= 3;
    return y;
}

int main(){
    int a = 1, b = 4;
    b = f(b, b);
    cout << a << ":" << b <<':';
    {
        int a = 7, b = 2;
        cout << ((b<=4) || ((a-=2) <= 4))
             << a << ":" << b << ":";
    }
    {
        inta = 8;
        cout << a << ':';
    }
    cout << a << ":" << b << endl;
    return 0;
}
%code

\par\vspace{5mm}
%enditem
%beginitem
9. %goodpath:Scripts/2018/Mod2018_1/01-good.txt
\textbf{Відмітити все, що є однією правильною лексемою:}
( 1 ) False

( 2 ) fop

( 3 ) <>

( 4 ) '123'
\par\vspace{5mm}
%enditem
%beginitem
10. \textbf{Що буде виведено за виконання фрагменту коду?}

%code
print('R', end=' ')
m=4
while m<=10:
    m+=2
print(m)
%code

\par\vspace{5mm}
%enditem










%end
\newpage

\newpage

%begin
{\centering\Large\bfseries Контрольна робота\par}
\medskip
\noindent\rule{\linewidth}{0.5pt}
\smallskip
\noindent\textbf{Варіант \No\,4}\hfill \textit{ПІБ:}\,\underline{\hspace{60mm}}
\vspace{4mm}

%beginitem
1. \textbf{Що буде виведено за виконання фрагменту коду?}
a = 4
b = 0
c = 2

def g(b):
global c    a -= 2
    b = 3
    c = 1
    return a + b + c

a = 5
b = 1
c = 6

try:
    print(g(b), a, b, c, sep=';')
except:
    print('Z', a, b, c, sep=';')
\par\vspace{5mm}
%enditem
%beginitem
2. 
Записати ВИРАЗ, значенням якого є множина, що складається з цілих чисел від 15 до 2005 (включно), що діляться на 3.\par\vspace{5mm}
%enditem
%beginitem
3. %goodpath:Scripts/2019/Mod2019_1/Lex/03-good.txt
\textbf{Відмітити все, що є коректним оператором С++:}
( 1 ) x=y=z

( 2 ) ++n

( 3 ) True=1

( 4 ) n=1; n++
\par\vspace{5mm}
%enditem
%beginitem
4. 
Записати ВИРАЗ, значенням якого є список, що складається з квадратних коренів цілих чисел від 15 до 2005 (включно), що діляться на 3, записаних в оберненому порядку.\par\vspace{5mm}
%enditem
%beginitem
5. %goodpath:Scripts/2018/Mod2018_1/01-good.txt
\textbf{Відмітити все, що є однією правильною лексемою:}
( 1 ) False

( 2 ) fop

( 3 ) <>

( 4 ) '123'
\par\vspace{5mm}
%enditem
%beginitem
6. 
Написати програму, що запитує у користувача значення дійсних параметрів $x$ та $y$, обчислює для них значення виразу 
$$\frac{\log_{10} x - 58}{(\arccos x - 0.6) (y - 52)}$$ 
та повiдомляє введенi данi та результат обчислення.
 Оцінюється правильність та якість коду.

\par\vspace{5mm}
%enditem
%beginitem
7. %goodpath:Scripts/2019/Mod2019_1/Lex/01-good.txt
\textbf{Відмітити все, що є однією правильною лексемою:}
( 1 ) 0.5

( 2 ) or

( 3 ) -3J

( 4 ) "abc"""
\par\vspace{5mm}
%enditem
%beginitem
8. \textbf{Що буде виведено за виконання фрагменту коду?}

print('R', end=' ')
for e in range(7,7,-2):
    if e<=7:
        continue
        print(e, end=' ')
        e=-5
    if e<=7:
        break
    else:
        print(e, end=' ')
    print(e)
\par\vspace{5mm}
%enditem
%beginitem
9. \textbf{Що буде виведено за виконання фрагменту коду?}
e=['0','1','2','3','4']; e.extend([1,2]); print(e)\par\vspace{5mm}
%enditem
%beginitem
10. 

\begin{center}
	\adjustimage{max size={0.8\linewidth}{0.8\paperheight}}
	{../15BinTree/trees_exp/40.png}
\end{center}
Указати послідовність міток вершин дерева, зображеного на рис., що утвориться в результаті обходу дерева в оберненому порядку. Мітки записувати через пробіл.\par Алгоритм починає роботу з кореня (на рис. це верхня вершина).
%answer:40 оберненому
\par\vspace{5mm}
%enditem










%end
\newpage
\end{document}


